\chapter{Accélérateurs pour le travail quantitatif}\label{ch:pl:accelerators-for-quantitative-work}

Une banque a transféré sa simulation de Monte Carlo du risque de contrepartie, exécutée chaque nuit, sur des processeurs graphiques, et elle a gagné un facteur de plusieurs fois. La même banque a essayé les
mêmes dispositifs pour le pricing intrajournalier de transactions isolées, et rien n'a bougé~: chaque requête passait plus de temps à traverser le lien vers le dispositif et à en revenir
qu'à calculer sur celui-ci. Les deux résultats étaient prévisibles à partir de trois nombres par charge de travail~--- le volume d'arithmétique effectué, le nombre d'octets
déplacés, et le nombre d'octets devant traverser le lien vers le dispositif~--- et deux nombres par machine. Ce chapitre mesure ces nombres sur le processeur du portable
et reprend celui d'un dispositif de centre de calcul depuis sa fiche produit, et prédit les deux issues avant que quiconque n'achète quoi que ce soit.

\section{Ce que fait bien un processeur graphique}

Un processeur graphique (un accélérateur au sens du livre~12, chapitre~23) dispose de milliers d'unités arithmétiques simples, qui exécutent la même
instruction sur des données différentes, et d'une mémoire qui lui est attachée, plusieurs fois plus rapide que celle d'un serveur. Il excelle dans ce que font déjà les
instructions SIMD d'un CPU (livre~13, chapitre~14), mais à une échelle bien plus grande~: une longue boucle d'opérations arithmétiques identiques et indépendantes sur des
tableaux déjà présents dans sa mémoire. Il est en revanche mal adapté à trois choses que le code quantitatif fait souvent.

\begin{definition}[Transfert hôte--périphérique]\label{def:pl:accelerators-for-quantitative-work:transfer}
Un \emph{transfert hôte--périphérique}\index{transfert hôte--périphérique} déplace des données entre la mémoire du CPU (l'hôte) et celle d'un accélérateur
(le périphérique) par un lien tel que PCI Express~; chaque entrée lue par un calcul sur le périphérique et chaque résultat qu'il renvoie traverse ce lien, à
une bande passante très inférieure à celle de l'une ou l'autre mémoire, et avec un coût fixe à chaque traversée.
\end{definition}

\begin{definition}[Divergence des threads]\label{def:pl:accelerators-for-quantitative-work:divergence}
La \emph{divergence des threads}\index{divergence des threads} se produit lorsque des threads qui s'exécutent en pas de deux prennent des branches différentes~: le groupe exécute
chaque branche à tour de rôle, les threads qui ne s'y trouvent pas restant inactifs, si bien qu'un code à branches dépendant des données (exercice anticipé, contrôles de barrière, gains dépendant du
chemin avec de nombreux cas) tourne à une fraction du débit du périphérique.
\end{definition}

\begin{omfigure}
\begin{tikzpicture}[font=\scriptsize,box/.style={draw,rounded corners=2pt,align=center,minimum height=0.8cm,inner sep=3pt},>=Stealth]
\node[box,fill=omMeth!15,minimum width=2.6cm] (cpu) at (0,0) {cœurs CPU};
\node[box,fill=gray!10,minimum width=2.6cm] (hm) at (0,-1.6) {mémoire hôte};
\node[box,fill=omDef!15,minimum width=3.2cm] (gpu) at (9.0,0) {périphérique~: des milliers\\ d'unités arithmétiques};
\node[box,fill=gray!10,minimum width=3.2cm] (dm) at (9.0,-1.6) {mémoire du périphérique};
\draw[<->,thick] (cpu) -- node[left]{24.0~Go/s, un cœur} (hm);
\draw[<->,very thick] (gpu) -- node[right]{3.35~To/s} (dm);
\draw[<->,thick,omThm] (hm.east) -- node[above]{PCIe Gen5~: 128~Go/s} node[below,align=center]{chaque entrée et\\chaque résultat traverse} (dm.west);
\draw[->,dashed] (cpu.east) -- node[above=2pt,align=center]{lancement et synchronisation~:\\$L$ par requête} (gpu.west);
\end{tikzpicture}

\omcaption{Où passe le temps lorsque le travail migre vers un périphérique. La mémoire du périphérique est 140 fois plus rapide que ce qu'atteint un cœur du portable,
mais les entrées et les résultats traversent un lien plus lent que l'une ou l'autre mémoire, et chaque requête paie un coût fixe de lancement et de synchronisation $L$
(le chapitre suppose 10~microsecondes).}\label{fig:pl:accelerators-for-quantitative-work:path}
\end{omfigure}

Le troisième est la taille. Un calcul qui prend une microseconde sur un CPU ne peut pas être plus rapide sur un périphérique dont le seul lancement et la seule synchronisation
prennent plusieurs microsecondes~: un petit travail, comme le prix d'une transaction isolée, supporte un coût fixe que les gros travaux amortissent.

\section{Le modèle roofline}

\begin{definition}[Intensité arithmétique]\label{def:pl:accelerators-for-quantitative-work:intensity}
L'\emph{intensité arithmétique}\index{intensité arithmétique} d'un noyau de calcul est le nombre d'opérations en virgule flottante qu'il effectue par octet
qu'il déplace entre la mémoire et le processeur.
\end{definition}

\begin{definition}[Modèle roofline]\label{def:pl:accelerators-for-quantitative-work:roofline}
Le \emph{modèle roofline}\index{modèle roofline} (Williams, Waterman et Patterson, 2009) borne le débit arithmétique qu'un noyau d'intensité
$I$ peut atteindre sur une machine de débit crête $P$ et de bande passante mémoire $B$ par $\min(P,\ I B)$. Tracé en axes logarithmiques, c'est une droite
oblique (la bande passante) rencontrant un palier horizontal (le débit crête) au point de jonction $I^* = P/B$.
\end{definition}

\begin{definition}[Noyau limité par la mémoire, noyau limité par le calcul]\label{def:pl:accelerators-for-quantitative-work:bound}
Un noyau dont l'intensité est inférieure au point de jonction de la machine est \emph{limité par la mémoire}\index{noyau limité par la mémoire} (memory-bound)~: son débit
est limité par la bande passante, et ajouter des unités arithmétiques ne l'aide pas. Un noyau situé au-dessus de la jonction est \emph{limité par le calcul}\index{noyau limité par le calcul} (compute-bound).
\end{definition}

Le roofline transforme une question sur le matériel en deux nombres sur le code. Le modèle du chapitre (\texttt{firm.roofline}) ajoute le transfert
et le coût fixe~: sur un périphérique, un noyau prend au moins $L + X/\ell + \max(F/P, M/B)$ secondes, avec $F$ ses opérations, $M$ ses octets
en mémoire, $X$ les octets qui traversent un lien de bande passante $\ell$, et $L$ la latence de lancement et de synchronisation d'une requête.

\omcode{code/firm/roofline/firm_roofline.py}{55}{75}{Un toit et le temps de bout en bout d'un noyau sur celui-ci~: la borne roofline, plus une
latence de lancement et un transfert par le lien hôte pour un périphérique.}\label{lst:pl:accelerators-for-quantitative-work:roof}

\section{Mesurer le toit du portable}

Un cœur de ce portable (un Intel Core Ultra 7 155H, sous WSL2) a été mesuré~: le débit crête par un produit de matrices $2000\times2000$ en double précision
sur un thread BLAS, la bande passante par un triad de type STREAM $a_i = b_i + s c_i$ en C++20 sur des tableaux de 128~Mo, en comptant 24 octets
par élément. Le cœur atteint 66.9~GFLOP/s et 24.0~Go/s~: son point de jonction se situe à 2.79 opérations par octet.

\omcode{code/firm/roofline/cpp/firm_roofline_triad.cpp}{11}{22}{La mesure de bande passante~: le triad, chronométré au meilleur de dix essais.}\label{lst:pl:accelerators-for-quantitative-work:triad}

Quatre noyaux de référence ont été chronométrés sur le même cœur (\cref{tab:pl:accelerators-for-quantitative-work:kernels})~: un pas de
mouvement brownien géométrique sur $2^{23}$ trajectoires ($S \mathrel{*}= e^{a + bZ}$), un gain de panier sur 500\,000 trajectoires de 20 actifs, l'agrégation d'un
cube de valeurs de transactions (2\,000 trajectoires, 50 dates, 200 transactions) en un profil d'exposition positive attendue, et le produit de matrices. Les
opérations sont comptées selon la convention de \texttt{firm.roofline}~: additions et multiplications comptent pour un, de même qu'un appel à \texttt{exp},
bien qu'il coûte de nombreux cycles.

\begin{table}[ht]
\centering\small
\begin{tabular}{lrrrr}
\toprule
noyau & intensité (flop/octet) & atteint (GFLOP/s) & toit (GFLOP/s) & part du toit \\
\midrule
pas de trajectoire & 0.17 & 0.65 & 3.99 & 16~\% \\
gain de panier & 0.25 & 4.82 & 5.99 & 81~\% \\
agrégation des expositions & 0.13 & 2.63 & 3.03 & 87~\% \\
produit de matrices & 166.7 & 64.2 & 66.9 & 96~\% \\
\bottomrule
\end{tabular}
\caption{Quatre noyaux sur un cœur du portable~: intensité selon la convention de comptage du chapitre, débit atteint (meilleur de cinq exécutions),
toit à cette intensité et part atteinte. Données~: \texttt{bench\_roofline.py}.}
\label{tab:pl:accelerators-for-quantitative-work:kernels}
\end{table}

Trois des quatre sont \omterm{def:pl:accelerators-for-quantitative-work:bound}{limités par la mémoire}, et les deux boucles simples atteignent 81 et 87~\% du toit de bande passante~: elles sont achevées, et seule
une bande passante supérieure les accélérerait. Le pas de trajectoire reste très en deçà de son toit parce que le roofline compte des opérations et que \texttt{exp} est une seule
opération qui prend de nombreux cycles~; sa vraie limite est le débit de l'exponentielle. Le produit de matrices, seul \omterm{def:pl:accelerators-for-quantitative-work:bound}{noyau limité par le calcul}, définit
le débit crête~: il atteint ici 96~\% de celui-ci seulement parce que le débit crête est son meilleur résultat sur davantage d'exécutions que les cinq du tableau.

\begin{omfigure}
\begin{tikzpicture}
\begin{loglogaxis}[width=12cm,height=6.4cm,xlabel={intensité arithmétique (flop par octet, échelle log)},ylabel={GFLOP/s (échelle log)},xmin=0.01,xmax=1000,
  ymin=0.1,ymax=100000,grid=major,grid style={gray!20},tick label style={font=\small},label style={font=\small},table/col sep=comma]
\addplot[omThm,very thick] table[x=intensity,y=cpu_gflops]{figdata/platforms/14-accelerators-for-quantitative-work/roofline.csv};
\addplot[omDef,very thick,dashed] table[x=intensity,y=device_gflops]{figdata/platforms/14-accelerators-for-quantitative-work/roofline.csv};
\addplot[only marks,mark=*,mark size=2pt,black]
  table[x=intensity,y=gflops]{figdata/platforms/14-accelerators-for-quantitative-work/measured_kernels.csv};
\node[font=\scriptsize,anchor=west] at (axis cs:0.18,0.65) {pas de trajectoire};
\node[font=\scriptsize,anchor=west] at (axis cs:0.29,4.8) {gain de panier};
\node[font=\scriptsize,anchor=south east,align=right,inner sep=1pt] at (axis cs:0.12,2.9) {agrégation\\des expositions};
\node[font=\scriptsize,anchor=south east] at (axis cs:167,64) {produit de matrices};
\node[omThm,font=\scriptsize,anchor=north] at (axis cs:30,50) {un cœur du portable};
\node[omDef,font=\scriptsize,anchor=north,align=center] at (axis cs:100,30000) {périphérique de centre\\de calcul (FP64)};
\end{loglogaxis}
\end{tikzpicture}

\omcaption{Deux rooflines~: un cœur du portable (mesuré) et un processeur graphique de centre de calcul en double précision (d'après sa fiche
produit), avec les quatre noyaux tels que mesurés sur le cœur. Le ridge du cœur est à 2.8 opérations par octet, celui du périphérique à 10.1~: les
noyaux \omterm{def:pl:accelerators-for-quantitative-work:bound}{limités par la mémoire} (memory-bound) sur le cœur le sont encore davantage sur le périphérique. Données~: \texttt{bench\_roofline.py}, \texttt{fig\_roofline.py}.}
\label{fig:pl:accelerators-for-quantitative-work:roofline}
\end{omfigure}

\begin{dated}{2026-09}{Un périphérique de centre de calcul}\label{dat:pl:accelerators-for-quantitative-work:h100}
La fiche produit de NVIDIA indique, pour le H100 dans sa forme SXM~: 34 téraFLOPS en double précision (67 sur ses cœurs tensoriels), 67 en simple
précision, 80~Go de mémoire à 3.35~To/s, une interconnexion NVLink à 900~Go/s, 128~Go/s par PCIe Gen5, et jusqu'à 700~W. Face à un cœur du portable, il offre 508
fois le pic en double précision et 140 fois la bande passante. Le chapitre prend ces chiffres comme le toit du périphérique et son
chiffre PCIe comme liaison avec l'hôte, la lecture la plus favorable de celui-ci.
\end{dated}

\section{Monte Carlo et risque sur un périphérique}

La charge de nuit est une simulation Monte Carlo d'expositions pour un ensemble de compensation (netting set)~: 10\,000 trajectoires de 50 facteurs de risque à
120 dates mensuelles, 5\,000 transactions valorisées sur chaque trajectoire et chaque date (dix opérations chacune~: une transaction linéaire sur un
facteur, actualisée), et les valeurs agrégées en un profil d'exposition attendue (livre~6, chapitres 17 et~18). Elle effectue 66~milliards d'opérations et déplace 98~Go en mémoire~: limitée par la mémoire sur les
deux machines. Seules les descriptions des transactions partent vers le périphérique (100 octets chacune) et seul le profil revient~; les trajectoires sont engendrées là où
elles servent, avec un générateur à compteur (livre~4, chapitre~26), de sorte que les trajectoires du périphérique sont celles du CPU.

\omcode{code/platforms/14-accelerators-for-quantitative-work/python/pl_roofline.py}{43}{65}{Les deux charges du modèle~: la tâche d'exposition
ses trois noyaux, son temps de bout en bout sur le périphérique, et une requête de valorisation de $n$ transactions.}\label{lst:pl:accelerators-for-quantitative-work:model}

Sur le cœur, à son toit, la tâche prend 4.09~secondes~; sur le périphérique, 29.3~millisecondes, lancements et transferts compris~: 140 fois plus
vite, soit le rapport des bandes passantes, car la tâche est limitée par la mémoire. Si 95~\% du traitement de nuit relève de ce type de calcul et que
le reste (lecture des positions, écriture des rapports) reste sur l'hôte, la loi d'Amdahl plafonne le gain de l'ensemble à $1/(0.05 + 0.95/140) = 17.6$ fois.
Face à un socket serveur à nombreux cœurs, le facteur par cœur se réduit~: un calcul \omterm{def:pl:accelerators-for-quantitative-work:bound}{limité par la mémoire} gagne le rapport de la bande passante du
périphérique à celle du socket, et non à celle d'un seul cœur.

\section{Entraînement, et ce que le livre d'apprentissage automatique disait déjà}

Le livre~12 (chapitre~23) a couvert le terrain propre de l'accélérateur~: l'entraînement des réseaux de neurones est dominé par les produits de matrices, \omterm{def:pl:accelerators-for-quantitative-work:bound}{limités par le calcul} (compute-bound), à
forte intensité, exécutés en précision mixte sur les unités tensorielles et répartis entre périphériques par parallélisme de données. Le même roofline l'explique~: un
produit de matrices de taille $n$ a une intensité de $n/12$, bien au-dessus de tout ridge, et le pic en précision réduite d'un périphérique est plusieurs fois son
pic en double précision. Le Monte Carlo de la finance quantitative est d'ordinaire l'autre type de noyau --- faible intensité, double précision --- et
il gagne le rapport des bandes passantes, non celui des pics.

\section{Décider avec un modèle de coût}

Une seule demande de prix intrajournalière est petite~: dix mille opérations (cent flux sur une courbe interpolée) et un kilo-octet traversant la
liaison (la transaction et son résultat). Sur le cœur, elle prend 0.33~microseconde~; sur le périphérique, 10.0 --- presque tout étant le lancement et
la latence de synchronisation, que le chapitre fixe à 10~microsecondes (une hypothèse, non un chiffre de fiche technique). Le périphérique
est 30 fois plus lent. Regroupées, les demandes partagent le coût fixe (\cref{fig:pl:accelerators-for-quantitative-work:batch})~: le périphérique atteint le seuil de rentabilité à 31 transactions par requête
et approche, pour de très grands lots, 32.7 fois le cœur --- un plafond imposé par la liaison, car chaque kilo-octet de transaction la traverse encore.

\begin{omfigure}
\begin{tikzpicture}
\begin{loglogaxis}[width=11.5cm,height=5.8cm,xlabel={transactions par requête (échelle log)},ylabel={accélération du périphérique (échelle log)},
  xmin=1,xmax=1000000,ymin=0.02,ymax=100,grid=major,grid style={gray!20},tick label style={font=\small},label style={font=\small},
  table/col sep=comma]
\addplot[omDef,very thick] table[x=n,y=speedup]{figdata/platforms/14-accelerators-for-quantitative-work/batch.csv};
\draw[black,dashed] (axis cs:1,1) -- (axis cs:1000000,1);
\draw[omThm,thick] (axis cs:31,0.02) -- (axis cs:31,100) node[pos=0.12,right,font=\scriptsize]{seuil de rentabilité~: 31 transactions};
\node[font=\scriptsize,anchor=north east] at (axis cs:900000,28) {plafond imposé par la liaison 32.7};
\end{loglogaxis}
\end{tikzpicture}

\omcaption{Accélération de bout en bout prédite du périphérique par rapport à un cœur du portable pour une requête valorisant $n$ transactions, avec une latence de lancement et
de synchronisation de 10~microsecondes et le débit PCIe de la fiche produit. Une seule transaction est 30 fois plus lente sur le périphérique~; le périphérique
l'emporte à partir de 31 transactions, et aucun lot ne le porte au-delà du plafond de la liaison. Données~: \texttt{fig\_roofline.py}.}
\label{fig:pl:accelerators-for-quantitative-work:batch}
\end{omfigure}

\begin{example}[Plus rapide la nuit, non en journée]\label{ex:pl:accelerators-for-quantitative-work:verdict}
Le même périphérique, le même style de code, deux verdicts. La tâche d'exposition est grande et limitée par la mémoire, elle déplace peu à travers la liaison et tourne 140
fois plus vite que sur un cœur (17.6 fois pour l'ensemble du traitement selon la loi d'Amdahl). La demande de prix est
minuscule, et la latence d'un lancement représente trente fois son coût total sur le CPU. Le lot de rentabilité dépend de la latence supposée~: 4 transactions à 1~microseconde, 16 à 5, 31 à 10,
155 à 50. Un service de valorisation capable de regrouper ses demandes à plusieurs centaines par appel gagnerait~; un service qui répond à chaque demande dès son arrivée ne le peut pas.
\end{example}

\begin{remark}[Ce que le modèle laisse de côté]\label{rem:pl:accelerators-for-quantitative-work:limits}
Le modèle est une borne~: les noyaux réels atteignent une fraction de l'un ou l'autre toit, comme le montre le tableau pour le cœur. Il ignore la \omterm{def:pl:accelerators-for-quantitative-work:divergence}{divergence des threads}, le coût de
portage et de la maintenance d'une seconde implémentation, le prix et la consommation du périphérique, et la mise en file d'attente des requêtes sur un périphérique partagé. Il suffit
à écarter un accélérateur, ou à dire que le mesurer vaut l'effort.
\end{remark}

\begin{dated}{2026-09}{Un indice de référence pour le matériel de valorisation et de risque}\label{dat:pl:accelerators-for-quantitative-work:stac}
STAC Research décrit sa suite de tests STAC-A2 comme «~la référence du secteur pour tester les piles technologiques utilisées pour les charges de calcul intensif
analytiques liées à la valorisation et à la gestion des risques~». Les fournisseurs publient des résultats sur celle-ci pour les CPU et les accélérateurs~; le
roofline du chapitre est l'estimation au dos d'enveloppe de la même question.
\end{dated}

\section{Tutoriel~: deux rooflines et une décision}\label{tut:pl:accelerators-for-quantitative-work:tut}

\begin{tutorial}
\textbf{Objectif.} Mesurer le toit d'un cœur et quatre noyaux sur celui-ci, les placer à côté du toit d'un périphérique cité, et prédire deux charges de
bout en bout. \textbf{État final~:} \cref{fig:pl:accelerators-for-quantitative-work:roofline,fig:pl:accelerators-for-quantitative-work:batch}.
\begin{tutsteps}
\item \textbf{Mesurer}~: \texttt{bench\_roofline.py} avec un seul thread BLAS~: le pic, la bande passante du triad, les quatre noyaux.
\item \textbf{Placer}~: \texttt{Kernel.intensity} et \texttt{Roof.attainable} pour chaque noyau sur les deux toits.
\item \textbf{Prédire}~: \texttt{overnight(cpu\_roof())}, \texttt{request(cpu\_roof(), n)}, \texttt{break\_even}.
\item \textbf{Éprouver l'hypothèse}~: \texttt{launch\_sensitivity}.
\end{tutsteps}
\textbf{Pour aller plus loin.} Mettre les chiffres en simple précision du périphérique dans le toit et voir quelles charges il modifie~; garder les données de marché
sur le périphérique entre les requêtes et n'envoyer que la transaction.
\end{tutorial}

\section{Construction~: le modèle roofline}\label{bld:pl:accelerators-for-quantitative-work:roofline}

\begin{build}
\bfield{Rôle} Un modèle de coût qui indique à la plateforme, avant tout achat ou portage, quelles charges un accélérateur accélérerait et de
combien, à partir de chiffres mesurés et cités.
\bfield{Interface} \texttt{Kernel(name, flops, bytes, transfer)}~; \texttt{Roof(name, peak, bandwidth, link, launch, source)} avec
\texttt{attainable}, \texttt{time}, \texttt{end\_to\_end}, \texttt{ridge}~; \texttt{DEVICES}~; \texttt{measure\_peak},
\texttt{measure\_bandwidth}~; les noyaux de référence et leurs descripteurs~; \texttt{amdahl}~; \texttt{break\_even}.
\bfield{Règles} Les chiffres des périphériques sont des données avec une source~; les chiffres du CPU sont mesurés sur un cœur avec un seul thread~; une convention de comptage
des opérations énoncée~; les hypothèses (la latence de lancement) signalées et balayées.
\bfield{Tests d'acceptation} \texttt{code/firm/roofline/tests/}~: l'arithmétique du toit~; chaque noyau de référence contre un calcul direct et son
intensité comptée~; la loi d'Amdahl~; la recherche du seuil de rentabilité, y compris un périphérique qui ne l'atteint jamais~; les données du périphérique et le triad.
\bfield{Pour aller plus loin} Un toit mesuré de tout le socket~; un toit de périphérique issu d'un benchmark plutôt que d'une fiche technique~; le coût par résultat en argent
à partir des prix de location.
\end{build}

\begin{omsources}
\item S.~Williams, A.~Waterman et D.~Patterson, ``Roofline: an insightful visual performance model for multicore architectures'',
\emph{Communications of the ACM} 52(4), 2009.
\item G.~M. Amdahl, ``Validity of the single processor approach to achieving large scale computing capabilities'', AFIPS Spring Joint
Computer Conference, 1967.
\item NVIDIA, page produit du H100~; STAC Research, page de l'indice de référence STAC-A2.
\end{omsources}

\section{Exercices}

\begin{exercise}[$\star$]\label{exo:pl:accelerators-for-quantitative-work:1}
Quelle est l'intensité d'un produit matriciel de taille $n$ selon la convention du chapitre, et à partir de quel $n$ est-il \omterm{def:pl:accelerators-for-quantitative-work:bound}{limité par le calcul} (compute-bound) sur le cœur et
sur le périphérique~?
\end{exercise}

\begin{exercise}[$\star$]\label{exo:pl:accelerators-for-quantitative-work:2}
Pourquoi l'accélération du travail d'exposition est-elle proche du rapport des bandes passantes et non du rapport des crêtes~?
\end{exercise}

\begin{exercise}[$\star$]\label{exo:pl:accelerators-for-quantitative-work:3}
Avec 95~\% d'un processus accéléré 140 fois, quel est le gain du processus entier, et quel serait-il avec 99~\%~?
\end{exercise}

\begin{exercise}[$\star\star$]\label{exo:pl:accelerators-for-quantitative-work:4}
Pourquoi l'étape de chemin n'atteint-elle que 16~\% de son toit, et qu'est-ce qui le rapprocherait de celui-ci~?
\end{exercise}

\begin{exercise}[$\star\star$]\label{exo:pl:accelerators-for-quantitative-work:5}
D'où vient le plafond de 32.7 pour les grands lots de requêtes de prix~?
\end{exercise}

\begin{exercise}[$\star\star$]\label{exo:pl:accelerators-for-quantitative-work:6}
Quelles parties d'une simulation de Monte Carlo d'option américaine avec une régression à chaque date souffriraient de la \omterm{def:pl:accelerators-for-quantitative-work:divergence}{divergence des threads}, et lesquelles non~?
\end{exercise}

\begin{exercise}[$\star\star\star$]\label{exo:pl:accelerators-for-quantitative-work:7}
\emph{Programmation.} Avec \texttt{firm.roofline}, recalculer le lot de seuil de rentabilité si les données de marché (1~Mo) doivent être envoyées avec chaque requête au lieu
de rester sur le périphérique.
\end{exercise}

\begin{exercise}[$\star\star\star$]\label{exo:pl:accelerators-for-quantitative-work:8}
\emph{Trouver la faille.} «~Le périphérique a 508 fois la crête de notre cœur, donc notre calcul de risque, qui prend dix heures sur un cœur, prendra une minute.~»
\end{exercise}

\section{Problème~: Plus rapide la nuit, pas en cours de journée}

\begin{problem}[{Problème du week-end --- deux charges de travail, un périphérique}]\label{pb:pl:accelerators-for-quantitative-work:1}
Un cœur de portable tel que mesuré, le périphérique de l'\cref{dat:pl:accelerators-for-quantitative-work:h100}, les deux charges de travail du chapitre.

\medskip\noindent\textbf{Partie I --- Toits.}
\begin{enumerate}
\item Quelles sont la crête, la bande passante et le point de bascule du cœur~?
\item Quelles sont celles du périphérique, et leurs rapports au cœur~?
\item Lesquels des quatre noyaux sont \omterm{def:pl:accelerators-for-quantitative-work:bound}{limités par la mémoire} (memory-bound) sur le cœur, et sur le périphérique~?
\item Pourquoi l'étape de chemin se situe-t-elle sous son toit~?
\item Que ne capture pas un roofline~?
\end{enumerate}

\medskip\noindent\textbf{Partie II --- Le travail de nuit.}
\begin{enumerate}[resume]
\item Combien d'opérations et d'octets le travail d'exposition comporte-t-il, et qu'est-ce qui traverse la liaison~?
\item Quels sont ses temps sur le cœur et sur le périphérique~?
\item Pourquoi l'accélération est-elle ce qu'elle est~?
\item Que donne la loi d'Amdahl pour le processus entier~?
\item Comment la comparaison changerait-elle face à un serveur à nombreux cœurs~?
\end{enumerate}

\medskip\noindent\textbf{Partie III --- La requête.}
\begin{enumerate}[resume]
\item Quels sont les temps d'une requête unique sur le cœur et sur le périphérique~?
\item À quel lot le périphérique atteint-il le seuil de rentabilité~?
\item Comment le seuil de rentabilité se déplace-t-il avec la latence de lancement~?
\item Qu'est-ce qui borne l'accélération des très grands lots~?
\item Que devrait changer un service de tarification pour en bénéficier~?
\end{enumerate}

\medskip\noindent\textbf{Partie IV --- Le verdict.}
\begin{enumerate}[resume]
\item Énoncer le \emph{résultat nommé}~: l'accélération de bout en bout prédite de la simulation de Monte Carlo d'exposition et d'une requête de prix portant sur une seule transaction, et
la taille du lot à partir de laquelle le périphérique atteint le seuil de rentabilité.
\item Quelles charges de travail de la plateforme faudrait-il déplacer d'abord vers un périphérique~?
\item Que faudrait-il mesurer avant d'en acheter un~?
\item Pourquoi utiliser la double précision dans le modèle~?
\item En une phrase~: quand un accélérateur est-il rentable~?
\end{enumerate}
\end{problem}

\section{Questions d'entretien}

\begin{interviewq}[$\star$ \iqroles{développeur}]\label{iq:pl:accelerators-for-quantitative-work:1}
Qu'est-ce que l'\omterm{def:pl:accelerators-for-quantitative-work:intensity}{intensité arithmétique}, et pourquoi compte-t-elle pour choisir le matériel~?
\end{interviewq}

\begin{interviewq}[$\star\star$ \iqroles{développeur, chercheur}]\label{iq:pl:accelerators-for-quantitative-work:2}
Votre simulation de Monte Carlo a été 30 fois plus rapide sur un processeur graphique, votre service de tarification a ralenti. Expliquez.
\end{interviewq}

\begin{interviewq}[$\star\star$ \iqroles{développeur}]\label{iq:pl:accelerators-for-quantitative-work:3}
Comment estimeriez-vous, sans le matériel, ce qu'un périphérique ferait pour une charge de travail~?
\end{interviewq}

\begin{interviewq}[$\star\star$ \iqroles{chercheur}]\label{iq:pl:accelerators-for-quantitative-work:4}
Pouvez-vous exécuter une simulation de Monte Carlo sur un périphérique et obtenir exactement les nombres du CPU~? Comment procéder~?
\end{interviewq}

\begin{interviewq}[$\star\star\star$ \iqroles{développeur}]\label{iq:pl:accelerators-for-quantitative-work:5}
Concevez le processus de décision pour déplacer des parties d'une plateforme de risque vers des accélérateurs.
\end{interviewq}
